\documentclass{article}
\usepackage[margin=3cm]{geometry}
\usepackage{amsmath,amssymb,mathtools}
\usepackage[noend]{algpseudocode}
\usepackage{capt-of}
\usepackage{hyperref}
\hypersetup{hidelinks}
\makeatletter
\providecommand*{\theHALG@line}{}
\renewcommand*{\theHALG@line}{\thefigure.\arabic{ALG@line}}
\makeatother
\pagestyle{plain}

% Objects and fields.
\newcommand{\genesis}{B_{\mathrm{gen}}}
\newcommand{\slot}{\operatorname{slot}}
\newcommand{\parent}{\mathsf{parent}}
\newcommand{\proposer}{\mathsf{proposer}}
\newcommand{\valindex}{\mathsf{val\_index}}
\newcommand{\round}{\mathsf{round}}
\newcommand{\rootf}{\mathsf{root}}
\newcommand{\committee}{\mathcal K}
\newcommand{\acp}{\mathsf{ac\_payload}}

% Chain state.
\newcommand{\targetParticipation}{\mathsf{target\_participation}}
\newcommand{\progressBit}{\mathsf{progress}}
\newcommand{\Qtarget}{Q_{\mathrm{target}}}
\newcommand{\Qprog}{Q_{\mathrm{prog}}}
\newcommand{\pendingFinality}{\mathsf{pending\_finality}}
\newcommand{\finalityParticipation}{\mathsf{finality\_participation}}
\newcommand{\Qfin}{Q_{\mathrm{fin}}}

% Store, frame and fork choice.
\newcommand{\T}{\mathcal T}
\newcommand{\hmax}{h_{\max}}
\newcommand{\AC}{\mathcal A}
\newcommand{\Fr}{\Phi}
\newcommand{\tree}{\mathsf{tree}}
\newcommand{\fcroot}{\mathsf{fc\_root}}
\newcommand{\allowed}{\mathsf{viable\_heads}}
\newcommand{\filtered}{\fname{viable\_tree}}
\newcommand{\gfPool}{\mathsf{gf\_votes}}
\newcommand{\attestations}{\mathsf{attestations}}
\newcommand{\fname}[1]{\text{\textsc{#1}}}
\newcommand{\latestConfirmed}{\mathsf{latest\_confirmed}}
\newcommand{\timestamp}{\mathsf{timestamp}}
\newcommand{\resolutionTime}{\mathsf{resolution\_time}}

% SA state.
\newcommand{\savedConfirmation}{\mathsf{saved\_confirmation}}
\newcommand{\lockEntries}{\mathsf{locks}}
\newcommand{\lockDeadline}{\mathsf{lock\_deadline}}
\newcommand{\restartTime}{\mathsf{restart\_time}}

\title{Goldilocks specification: relation to the Lean model and claims}
\author{}
\date{}

\begin{document}
\maketitle

This note is a companion to \texttt{goldilocks-spec.tex}.  It lists the
differences between that specification and the Lean model, and the
continuation claims that are proved for the Lean model.  Its section
numbers refer to this note; the names of functions and fields are those
of the specification.  The specification is titled Goldilocks, and it uses that name
also for its available chain; this note calls the available chain SA (the
old Lean model's names use FGSA for the protocol).

\section{Relation to the Lean model}\label{sec:lean}

The Lean model follows \texttt{goldilocks-round-model.tex}.  It keeps two stores,
one for the FG and one for SA, and a per-block lock map.  This section
lists the differences.  Items (a)--(g) change no output and no message
of an honest node.  Items (h)--(k) change the behaviour.
Section~\ref{sec:fast-path} lists the later changes of the specification
that the Lean model does not have, loose BJ (P8) and the weak E1 (P9)
among them.

\paragraph{Locks.}
(a) The non-genesis blocks that pass the fast or slow test at a write form
a prefix of the path from $\Sigma.F$ to the result in $\Fr.\tree$: the score cannot decrease toward ancestors,
and a saved confirmation that extends~$B$ extends its ancestors.
(b) The live entries lie on one chain, because the veto rejects a result
that conflicts with any of them.
(c) For each non-genesis~$B$ that extends~$\Sigma.F$, the per-block
deadline is the maximum deadline over surviving writes whose top is at or
above~$B$; this is exactly the maximum deadline of the live entries at
or above~$B$ (a value that this form does not compute), including after expiry
and after the finality cleanup of the per-block form.  (That cleanup
deletes only deadlines of blocks that conflict with~$\Sigma.F$, and no
block at or above~$\Sigma.F$ is below such a block.)
(d) Release of~$B$ in the Lean model removes every block above~$B$ too, by opponent
monotonicity, so keeping the surviving prefix is exactly the move-back rule
of $\fname{unlock}$.  When $\Sigma.F=\genesis$, $\fname{unlock}$ can move
a lock to~$\genesis$, where the Lean model keeps no deadline.  Such a lock
has no effect, as a lock on~$\Fr.F$ has none: it vetoes no result and does
not change the anchor.
(e) The per-block form also keeps deadlines on strict ancestors
of~$\Sigma.F$, deletes the deadline of every block that conflicts with a
new~$\Sigma.F$, and replaces each saved confirmation that conflicts
with~$\Sigma.F$ by~$\bot$.  This form writes no lock below~$\Sigma.F$,
reads no content and no ancestry of a block outside $\Fr.\tree$ (it reads
the deadline of every entry, and it tests whether the block of an entry or
of a saved confirmation is in the tree), and deletes
nothing when $\Sigma.F$ advances.  No difference is observable.  A lock
on $B\prec\Sigma.F$ is outside every viable tree, so it is never the
anchor and never $\mathit{fg\_confirmed}$; it vetoes no result, because every
result extends~$\Sigma.F$; and it has no opponents, because the head of a
resolved vote extends~$\Sigma.F$.  A lock on a block that conflicts
with~$\Sigma.F$ locks no block that extends~$\Sigma.F$, and $\Sigma.F$
never moves back.  A saved confirmation outside $\Fr.\tree$ extends no
block of $\Fr.\tree$, so it fails the slow test for every block that this
form can lock.

\paragraph{Votes and state.}
(f) The Lean model resolves a vote when its head is held and compatible
with~$\Sigma.F$; this form asks for a head in $\Fr.\tree$.  The
difference is a vote whose head is a strict ancestor of~$\Sigma.F$.  Such
a vote supports no block of $\Fr.\tree$ and opposes no block of a locked
chain, so no score, walk, confirmation or unlocking changes at this node.
Item~(h) gives the one effect on a message.
(g) The Lean model keeps the fast-confirmed set as a store field
(\texttt{fastWrite}).  This form has no such field: the fast bit of a lock
marks the lock of a fast or optimistic confirmation (item~P5 of
Section~\ref{sec:fast-path}).  The block $\mathit{fg\_confirmed}$ of this form is
the \emph{source} of the Lean model (\texttt{actionSourceFG} and the
source names of the claims); the name changed to avoid the FFG meaning of
source.

\paragraph{Other representation.}
The Lean model writes the record of the confirmed tip as $\Sigma.F$ when
it does not extend~$\Sigma.F$; this form tests that in
$\Call{ac\_confirmed}{}$.  This form also cuts the output to the head:
$\Call{ac\_confirmed}{}$ returns the deepest common ancestor of that block
and the head of $\Call{ac\_head}{}$.  The record itself, and every rule that
reads it, are unchanged; the Lean model and its claims read the record,
not the cut output.  $\Call{ac\_head}{}$ is a new output with no Lean
counterpart: the fork choice at read time, with every held vote of the
latest slot whose votes are due (slot~$s-1$ before $t_s+2\Delta$, slot~$s$
from then on), on the current frame and locks.  It saves no head.
This form has continuous time and stamps an object with its receipt
time; the Lean model stamps it with the time of the last tick
(\texttt{State.now}).  They agree while the node is awake: every stamp
test is a strict comparison with a public time~$c$, all public times are
multiples of~$\Delta$, and an awake node runs a tick at each of them, so
the receipt time is below~$c$ exactly when the last tick time is.  An
asleep node runs no tick (item~(k)), so the agreement needs one more rule
in the Lean port: a restart sets \texttt{now} to the exact wake time~$w$.
Then an object received at a time $x$ before the next tick has the stamp
$w$, and no public time lies in $(w,x]$, so $w<c$ exactly when $x<c$.
Without that rule, the objects delivered at the wake would get the stamp
of the last tick before the sleep and could count as early support.  The
same exact time is $\AC.\restartTime$; a time rounded down to the last
tick can pass the join test $\AC.\restartTime\leq t_{\slot(t)-1}$ wrongly.
The Lean model keeps the explicit cutoff \texttt{lateVotes} (stamp below
the gate time) in the confirmation update; this form reads the whole
pool, which is the same set because the tick runs before the objects of
its time.  The Lean model uses the witness height
$\Sigma.\hmax-1$ in finality acceptance, also in a contest; this form uses
the head height floor.  They agree: outside a contest the two heights are equal,
and in a contest the block that carried the justification of~$\Sigma.J$
has state height $\Sigma.h_j+1=\Sigma.\hmax$ and extends every candidate,
because a candidate is an ancestor of~$\Sigma.J$.
\texttt{fgsa-spec-delta.md} lists the remaining differences of form.

\paragraph{Changes of behaviour.}
(h) \emph{The support list.}  By~(f), a proposer omits from
$B.\mathsf{support\_votes}$ a previous-slot vote whose head is a strict
ancestor of its~$\Sigma.F$.  The Lean model lists it.  The block content
and the block root differ, and a voter whose~$\Sigma.F$ is below that
head can count one support entry less.
(i) \emph{The contest.}  The Lean model tests the contest on every
accepted block; this form tests it on the blocks that
extend~$\Sigma.F$.  The two tests differ only when a block that conflicts
with~$\Sigma.F$ has a state height above the height at which $\Sigma.F$
was finalized.  That needs a quorum of finality pairs for $\Sigma.F$ and
a quorum of height pairs for another entry at the same height, so at
least $2q-n$ validators signed both.
(j) \emph{The vote merge.}  The Lean model adds to the frozen sets of a
voter the votes and the support entries of every accepted slot-$s$
proposal.  This form adds only those of the proposals in $\Fr.\tree$.
The two differ when a slot-$s$ proposal was accepted and $\Sigma.F$ then
advanced, in the same slot, to a block that conflicts with it.  The voter
cannot vote for that proposal; the Lean model still counts the votes that
it carries, and this form does not.
(k) \emph{Sleep and restart.}  In the Lean model an asleep node still
receives objects and runs its ticks and handlers; it only runs no duty
and sends nothing.  A Lean restart keeps only~$\Lambda$ and the clock: since
ruling~Y (on main from 6c36d481), \texttt{now} keeps the time of the last
tick.  The sleep port replaces that with the exact wake time (see the stamp
rule above).  In this form an asleep node receives nothing and runs no tick, and
a node that wakes restarts.  The restart keeps~$\Lambda$, $\Sigma$, the
clock, the locks and the record, and it resets the rest of~$\AC$.  A
restarted node unlocks, votes, proposes, runs the confirmation update and
gives a block to the FG only when it is joined at the time~$t$ of the
action ($\AC.\restartTime\leq t_{\slot(t)-1}$, one uniform test): a
restart erases votes that the network does not deliver again.  A node
that runs while its delivery is under the control of the adversary is
faulty.  The
Lean claims assume no restart and use the Lean sleep model, so they do not
cover this item.  The
port is planned in \texttt{drafts/sleep-join-plan.md}: one premise~$E$ in
the shape of RLMD-GHOST Definition~3, with the locks in place of the
retained votes.

The claims of Section~\ref{sec:continuation-claims} are proved for the
Lean model.  They do not cover items (h)--(k).

\section{The fast path: not yet in the Lean model}\label{sec:fast-path}

The specification (research commit \texttt{57ce025} and later) has these
changes that the Lean model of record does not have.  That model is the
repository \texttt{\textasciitilde/ethereum/lean/available-chain}, on
\texttt{main} at \texttt{d72d3bc3}: the round model of the earlier
specification.  The items of Section~\ref{sec:lean} and the claims of
Section~\ref{sec:continuation-claims} refer to it, and none of those
claims covers these changes.
\begin{itemize}
\item[(P1)] \emph{Fast confirmation.}  Unlocking counts equivocators as
  opponents and falls back to~$\Fr.F$; the locks of slot~$s$ and of later
  slots keep their place in the slot-$s$ unlocking.  The optimistic
  confirmation of slot~$s$, before $t_s+2\Delta$, locks the deepest block
  whose slot-$s$ votes with heads that extend it are three quarters of the
  committee, equivocators included, and that conflicts with no lock.  It
  has no restriction to blocks of slot~$s$ and no once-per-slot marker.
\item[(P2)] \emph{The FG update without rounds.}  After each tick the node
  asks the AC for a new pair; the FG client record decides whether an
  attestation follows.
\item[(P3)] \emph{Responsive voting.}  A node votes at the first ready
  proposal in $[t_s,t_s+\Delta]$, after an unlocking on the merged view.
  The proposer unlocks before it proposes, on its whole pool, and again
  before its vote; the second call can change the locks, because the own
  block can change the frame.  Receipt handlers can send votes, and the
  own vote is a receipt.
\item[(P4)] \emph{Interface.}  The receipt handlers take the frame.
\item[(P5)] \emph{Locks.}  A lock is $(B,j,f,e)$: a block, the slot~$j$ of
  the confirmation that wrote it, the fast bit~$f$, and its expiry~$e$,
  fixed at its write with $j+2\leq e\leq j+k$.  It is live while
  $\slot(t)\leq e$ and $B\in\Fr.\tree$.  A lock that is not fast gets
  $e=j+k$; a fast lock gets $\fname{fast\_expiry}$, the latest expiry of
  the node's live fast locks plus~$c\geq2$, at least $j+2$ and at most
  $j+k$, except that the update's fast lock keeps $j+k$ when a block of
  the result's chain passes the slow test.  The move back, unlocking and
  a restart keep the expiry.
  Every rule reads the locks through $\fname{live\_locks}$, which moves a
  lock back to its deepest ancestor that is viable or at or below
  $\Fr.\fcroot$.  Every lock write and every write of the record store
  the live locks, the new locks included ($\fname{add\_locks}$, also with
  no new lock), so a new lock on a nonviable block is stored on its viable
  ancestor.  A
  lock of slot~$s$ with the expiry $s+k$ expires at $t_{s+k+1}$, so a
  regular such lock lives $(4k-2)\Delta$ after its write, not
  $(4k-1)\Delta$; no lock lives longer, and there is no expiry guard.  The confirmation result is the deepest block at or
  above~$\Fr.F$ with a strict majority, with no walk.  An accepted update writes a fast lock on its deepest fast block
  and a lock that is not fast on a deeper slow block.
\item[(P6)] \emph{The FG source.}  $\fname{ac\_to\_fg}$ returns the
  deepest viable block under a live fast lock of slot $\slot(t)-2$ or
  later, when it is new.  There is no fast-confirmed field and no
  optimistic-confirmed field.
\item[(P7)] \emph{The confirmed tip} is the deepest common ancestor of the
  record and the head.
\item[(P8)] \emph{Loose BJ} (\texttt{shared/fg.tex}).  A justification at
  a height whose previous height is justified waits for the finality of its predecessor.  The client
  reads its finality pair from the entry state of its source, not from the
  head: at such a height the predecessor whenever the record allows it,
  with no test of finalized status, and at another height the newest
  pending pair that the record allows; there is no $h\leq h_f$ clear.
  $K_W=2$, and $\delta_{\mathrm t}\geq x+c_F+k_F$ in place of
  $\delta_{\mathrm t}=2R$.  The Lean transition justifies at every target
  quorum, its consumption invariant (I7) says so, and its window test
  \texttt{progReady} can cross a blocked target quorum before the timeout.
  None of the claims of Section~\ref{sec:continuation-claims}, and no Lean
  finality claim, covers this item.
\item[(P9)] \emph{Weak E1} (\texttt{shared/fg.tex}).  E1 needs the height
  pair at~$h$ to name a block other than the block of the finality pair
  at~$h$, for a timeout as for a target vote: a timeout $(h,T,\mathsf{true})$
  and a finality pair $(h,T)$ are not slashable.  The client record keeps
  at $\mathsf{timeout}[h]$ the set of the blocks that its timeouts at~$h$
  name, $\Call{permits}{}$ allows $(j,T)$ when that set is within~$\{T\}$,
  and $\Call{height\_pair}{}$ tests the lock before the timeout.  The Lean
  evidence predicate (\nolinkurl{legacy/AvailableChainSimplexFG/Evidence.lean}:13--17,
  \texttt{conflictsWithFinality}) makes every timeout at~$h$ conflict
  with every finality pair at~$h$; it needs the root test on its timeout
  arm.  The Lean client record keeps a Boolean timeout, and its client
  safety reads the old order of $\Call{height\_pair}{}$
  (\texttt{ClientSafety.lean}).  No Lean claim covers this item.
\end{itemize}
The port is planned in \nolinkurl{codex/jobs/fastpath-lean-plan-v2.md}
(v6, written for \texttt{332213d}) and the G2--G4 plans of the parent
folder: a separate exact evaluator, so that a vote can be sent at a
receipt; a new protocol instance for the specification (the eager
instance), next to the round instance; receipt tests at the receipt time.

\paragraph{Plan items that are stale for 57ce025.}
The proof map \nolinkurl{codex/reports/ac-proofmap-astra-p1.md} found
these.  The planned equality of two optimistic confirmations of one slot
is false: a node can confirm an ancestor and then a descendant, so the
claim is compatibility.  The plans still select the source from the latest
fast write and the optimistic field, use the expiry-guard equality and the
old deadline, and suppress the proposer's second unlocking.  The release
count at~$\Fr.F$ is the number of equivocators, not zero.

\paragraph{New contracts.}
\begin{itemize}
\item \emph{Lock to tip.}  A block under a live lock is at or below the
  effective record, and so at or below the confirmed tip.  This needs P5:
  without the normalization at every record write, a lock that reads as
  moved back can read as locked again above a record on a conflicting
  branch.
\item \emph{Optimistic protection.}  The note ``Why an optimistic
  confirmation is a fast confirmation'' gives the proof shape (the first
  conflicting honest vote).  The protection starts at
  $t_{s+1}+6\Delta$ for an optimistic confirmation of slot~$s$, as for a
  fast confirmation of slot~$s$; the claims name the slot of origin.
\item \emph{Recent source, next vote.}  When $\fname{ac\_to\_fg}$ returns
  $B$ in slot~$m$ from a fast lock of slot~$j$, $j\leq m\leq j+2$, every
  honest joined vote of slot~$m+1$ extends~$B$, under stated delivery,
  acceptance, retention, restart, viability and participation
  conditions.  For $k=2$ the source lock expires at~$t_{j+3}$, so the last
  step needs transferred votes, a new lock or the fork-choice root.  The
  outage continuation claim does not give this: it gives compatibility,
  not extension, and honest heads can move back.  A paper proof by
  transferred honest support, with no renewed lock, is in
  \nolinkurl{codex/reports/recent-source-astra-q1.md}, with named
  premises (source service from slot~$j$, batch service, honest votes of
  half of the committee in each slot, the voting path).  It can reuse the
  patterns of \texttt{healing\_first\_transfer\_fast} and
  \texttt{healing\_next\_transfer\_core}; the optimistic first transfer is
  new.
\item \emph{Source age.}  From slot $\slot(\tau)+3$ on, every source is a
  fast lock whose confirmation was written after~$\tau$.  The model must
  keep the origin slot of a lock through unlocking, the move back and
  restart; the lock stores it.
\end{itemize}
Preservation of the finality deadlines with attestations after every tick
is an open proof obligation.

\section{Continuation claims}\label{sec:continuation-claims}

Write $S_x(t)$ for validator~$x$'s inclusive available-output read.
Both public outputs, stable and confirmed, equal~$S$.  Write $A\parallel B$
for prefix compatibility, and put $L=(4k-1)\Delta$.

\paragraph{Outage continuation.}
Keep the healthy prefix and the fast or slow lag of $\textsf{Outage}_2$.
For $B\preceq S_x(T)$ at an honest reader, replace its old budget by
\[
 b_1-b_0+\bigl((b_0-6\Delta)\bmod 4\Delta\bigr)\leq L. \tag{V}
\]
Assume full-run synchrony with acceptance relay, $\mathrm{GST}\leq b_1$,
and $h_t>a$ for $b_1\leq t\leq\mathrm{horizon}$.  Participation is not
required in $[b_0,b_1)$.  There is no \textsf{FinalitySlashableBound}
premise.  For every honest~$y$ and $b_0\leq t\leq\mathrm{horizon}$,
\[
 B\preceq S_y(t),\qquad B\preceq\mathrm{stable}_y(t),\qquad
 B\preceq\mathrm{confirmed}_y(t),\qquad B\parallel F_y(t).
\]
\textsf{OutageContinuationClaim} is proved in Lean on \texttt{main},
commit \texttt{a3983c38}.  The old budget
$b_1-b_0\leq\min(\eta,L)$ does not suffice for this continuation:
\newline\texttt{AvailableChainFGSAProofs/Tests/OutageAfterWindowCounterexample.lean}\newline
on \texttt{main} gives a counterexample.  Sources: \nolinkurl{drafts/outage-output-direct.md},
\S\S1,10, and \nolinkurl{drafts/outage-persist-reserve.md}.

\paragraph{Weak continuation and corruption growth.}
Set $\delta_{\mathrm t}=2R$.  A strong healing prefix with
$a_{\mathrm{source}}\leq\lfloor(n-1)/3\rfloor$ selects one gate~$d_u$
and seed~$Q$ before the continuation is chosen.  The prefix ends at
$T_F=\mathrm{GST}+B_H^F$.  The ordinary weak continuation keeps the
same honest set and the complete event prefix.  Its safety, persistent
$6\Delta$ proposal inclusion, and growth use the same selected cut.
For every continuation gap $G'\geq4\Delta$, awake honest proposer
recurrence from~$d_u$ gives joint growth of stable and confirmed by
$t+G'+6\Delta$, when that deadline is within the horizon.

The corruption-growth companion permits a retained honest set
$H_{\mathrm{ret}}\subseteq H_{\mathrm{source}}$.  Each retained
validator has the same ordered local events, inputs, and awake values
through the inclusive endpoint~$T_F$.  Keep full-run synchrony,
acceptance relay, and original pending receipt deadlines.  Require
$2a<n$ at handover and $a<h_t$ after~$T_F$, so corruption can reach
$\lfloor(n-1)/2\rfloor$.  Keep \textsf{FinalitySlashableBound} for
all relevant carrier chains, including historical rows.  The companion
gives the same safety, inclusion, and conditional joint growth for
retained honest validators.  Its event index is computed in each run;
the gate and seed stay fixed.  \textsf{WeakContinuationGrowthClaim}
is proved in Lean on \texttt{main}, commit \texttt{8b09725a}.
\texttt{nvGrowth\_complete} is the non-vacuity witness.  Sources: \nolinkurl{drafts/wc-corruption-growth.md}
and \nolinkurl{drafts/wc-slashing-necessity.md}.

\paragraph{Proof errata.}
The P3 source-FG frame statement requires
$c_0\leq\mathrm{ht}(\genesis)\Rightarrow\Phi(\genesis)$.
Certified entries include genesis and all crossed entries.
Frame conclusions use the store and exact-FG frame premises and cover
every contest witness.  The tail proof also covers mixed finality
states: relay the finalized head and its certificate rows to the
terminal proposer, then use window retention and row selection.
K-TAIL r3 closes this case in Lean through \textsf{TailKernelClaim\_W}.
Details: \nolinkurl{../blueprint/22-finality.md},
\nolinkurl{../blueprint/23-continuation.md}, and the complete protocol.

\end{document}
